% PCCSL 507 Machine Learning Lab - Experiment 3 (pandas)
\documentclass[11pt,twoside]{article}

\usepackage[paper=letterpaper,top=0.7in,bottom=0.63in,left=0.75in,right=0.75in,headheight=24pt,headsep=0.14in,footskip=0.38in]{geometry}

\usepackage[T1]{fontenc}
\usepackage{mathptmx} % Times clone
\usepackage{xcolor}
\definecolor{primary}{HTML}{1F3864}
\usepackage{fancyhdr}
\usepackage{array}
\usepackage{colortbl}
\usepackage{parskip}
\setlength{\parindent}{0pt}
\setlength{\parskip}{4pt}
\setlength{\abovedisplayskip}{4pt}
\setlength{\belowdisplayskip}{4pt}
\setlength{\abovedisplayshortskip}{2pt}
\setlength{\belowdisplayshortskip}{2pt}
\usepackage{graphicx}
\usepackage{amsmath}
\usepackage{listings}
\lstset{basicstyle=\footnotesize\ttfamily, breaklines=true, frame=single, showstringspaces=false}

\pagestyle{fancy}
\fancyhf{}
\fancyhead[C]{\textcolor{primary}{\textbf{\small PCCSL 507 Machine Learning Laboratory Record}}}
\renewcommand{\headrulewidth}{0.75pt}
\renewcommand{\headrule}{\hbox to\headwidth{\color{primary}\leaders\hrule height \headrulewidth\hfill}}
\fancyfoot[C]{\footnotesize DEPARTMENT OF CSE \hfill CAPE COLLEGE OF ENGINEERING ALAPPUZHA \hfill Page No. \thepage}
\renewcommand{\footrulewidth}{0pt}

\newcommand{\expheading}[1]{{\fontsize{14}{17}\selectfont\textbf{#1}\par}}
\newcommand{\fullrule}{\noindent\rule{\textwidth}{0.6pt}\par}

\begin{document}
\sloppy

% Sheet-local tightening so it fits one right-side page
{\setlength{\parskip}{1pt}\setlength{\abovedisplayskip}{2pt}\setlength{\belowdisplayskip}{2pt}
\begin{center}
\expheading{EXPERIMENT NO. 3}
\end{center}
\vspace{2pt}

\noindent\textbf{Date:} \underline{\hspace{2.6cm}} \hfill \textbf{Page No.:} \underline{\hspace{2.6cm}}\par
\vspace{6pt}

\expheading{TITLE}
\vspace{4pt}
\noindent Pandas DataFrame -- Student Details\par
\vspace{4pt}

\expheading{OBJECTIVE}
\vspace{2pt}
\fullrule
\vspace{4pt}
\noindent To work with a Pandas DataFrame of five students (roll no, name, department, marks) and perform selection, filtering, aggregation and CSV export operations. \par
\vspace{4pt}

\expheading{AIM}
\vspace{4pt}
\noindent To create a student dataframe, display rows/columns, filter rows, compute summary statistics with department-wise average, and save it as CSV. \par
\vspace{4pt}

\expheading{THEORY}
\vspace{4pt}
DataFrame with rows and named columns:
\[\text{df}=\text{pd.DataFrame}(\{\text{columns}\})\qquad\text{df.head}(n),\; \text{df.tail}(n)\]
Column selection and boolean filtering:
\[\text{df}[[\text{``Name''},\text{``Marks''}]]\qquad \text{df}[\text{df}[\text{``Marks''}]>80]\]
Summary statistics:
\[\max,\; \min,\; \text{mean}=\frac{1}{n}\sum x_i,\; \text{sum}=\sum x_i\]
Department-wise average and CSV export:
\[\text{df.groupby}(\text{``Department''})[\text{``Marks''}].\text{mean}()\qquad \text{df.to\_csv}(\text{file})\]
\vspace{4pt}

\expheading{PROCEDURE}
\vspace{4pt}
\noindent 1. Create a DataFrame of 5 students with Roll no, Name, Department, Marks and display it.\par
\noindent 2. Display the first 3 rows (head), last 2 rows (tail) and only Name/Marks columns.\par
\noindent 3. Filter students with marks above 80 and students of CSE department.\par
\noindent 4. Compute highest, lowest, average and total marks; average mark department-wise.\par
\noindent 5. Save the DataFrame to CSV and verify the file. \par
\vspace{4pt}

\expheading{DATASET DESCRIPTION}
\vspace{4pt}
\noindent Dataset Name: Student details (entered in program)\par
\noindent Source: 5 students; departments CSE, ECE, ME\par
\noindent Number of Samples: 5 rows \hspace{0.3cm} Number of Features: 4 columns\par
\noindent Target / Class: Not applicable (data handling operations)\par
\noindent Description: Marks range 68--90; exported to exp3\_students.csv. \par
\vspace{4pt}
} % end sheet-local tightening

\newpage
\expheading{PROGRAM}
\vspace{6pt}

\begin{lstlisting}
import pandas as pd


data = {
    "Roll no": [1, 2, 3, 4, 5],
    "Name": ["Anu", "Ravi", "Meera", "Arjun", "Divya"],
    "Department": ["CSE", "ECE", "CSE", "ME", "CSE"],
    "Marks": [85, 72, 90, 68, 78],
}
df = pd.DataFrame(data)
print("Dataframe:")
print(df.to_string(index=False))

print("\nFirst 3 rows:")
print(df.head(3).to_string(index=False))
print("\nLast two rows:")
print(df.tail(2).to_string(index=False))
print("\nName and Marks columns:")
print(df[["Name", "Marks"]].to_string(index=False))
print("\nMarks more than 80:")
print(df[df["Marks"] > 80].to_string(index=False))
print("\nCSE department:")
print(df[df["Department"] == "CSE"].to_string(index=False))
print("\nHighest =", df["Marks"].max(), " Lowest =", df["Marks"].min(),
      " Avg =", round(df["Marks"].mean(), 2), " Total =", df["Marks"].sum())
print("\nAverage mark department wise:")
print(df.groupby("Department")["Marks"].mean().round(2).to_string())

df.to_csv("exp3_students.csv", index=False)
print("\nSaved to exp3_students.csv")
\end{lstlisting}

\newpage
\expheading{OUTPUT}
\vspace{2pt}
\noindent\rule{\textwidth}{0.5pt}\par
\vspace{4pt}
{\scriptsize
\begin{verbatim}
Dataframe:
 Roll no  Name Department  Marks
       1   Anu        CSE     85
       2  Ravi        ECE     72
       3 Meera        CSE     90
       4 Arjun         ME     68
       5 Divya        CSE     78

First 3 rows:
 Roll no  Name Department  Marks
       1   Anu        CSE     85
       2  Ravi        ECE     72
       3 Meera        CSE     90

Last two rows:
 Roll no  Name Department  Marks
       4 Arjun         ME     68
       5 Divya        CSE     78

Name and Marks columns:
 Name  Marks
  Anu     85
 Ravi     72
Meera     90
Arjun     68
Divya     78

Marks more than 80:
 Roll no  Name Department  Marks
       1   Anu        CSE     85
       3 Meera        CSE     90

CSE department:
 Roll no  Name Department  Marks
       1   Anu        CSE     85
       3 Meera        CSE     90
       5 Divya        CSE     78

Highest = 90  Lowest = 68  Avg = 78.6  Total = 393

Average mark department wise:
Department
CSE    84.33
ECE    72.00
ME     68.00

Saved to exp3_students.csv
\end{verbatim}}
\vspace{4pt}

\expheading{RESULT}
\vspace{4pt}
{\fontsize{12}{15}\selectfont
The above program was implemented and executed successfully using Python, and the corresponding output was obtained and verified. The results of the \textbf{Pandas DataFrame operations} were analyzed and found to be satisfactory, thereby achieving the desired objective of the experiment.\par
}

\vspace{12pt}

\noindent\textbf{Faculty Signature:} \underline{\hspace{3.6cm}} \hfill \textbf{Date:} \underline{\hspace{2.8cm}}\par

\end{document}
